
%%%%%%%%%%%%%%%%%%%%%%%%%%%%%%%%%%%%%%%%%%%%%%%%%%%%%%%%%%%%%%%%%%%%%%%%%%%%%%
%
% 11_adjunction.tex
%
%%%%%%%%%%%%%%%%%%%%%%%%%%%%%%%%%%%%%%%%%%%%%%%%%%%%%%%%%%%%%%%%%%%%%%%%%%%%%%

\begin{mosbox}{11. The Lift--Project Adjunction}

\BodyFont

MOS distinguishes two complementary representations of reasoning.

Natural language serves as the interface through which problems are
posed, while internally the system operates on validated reasoning DAGs.

Two boundary morphisms connect these worlds:

\[
\boxed{
\pi : \text{Language}
\longrightarrow
\text{Reasoning DAG}
}
\]

and

\[
\boxed{
G :
\text{Reasoning DAG}
\longrightarrow
\text{Language}.
}
\]

The lift converts a natural-language request into a validated reasoning
program, while the projection renders a completed reasoning graph back
into human-readable form.

\end{mosbox}

%%%%%%%%%%%%%%%%%%%%%%%%%%%%%%%%%%%%%%%%%%%%%%%%%%%%%%%%%%%%%%%%%%%%%%%%%%%%%%
%% HERO DIAGRAM
%%%%%%%%%%%%%%%%%%%%%%%%%%%%%%%%%%%%%%%%%%%%%%%%%%%%%%%%%%%%%%%%%%%%%%%%%%%%%%

\begin{center}

\begin{tikzpicture}[
  scale=1.0,
  dag node/.style={circle, fill=MOSRed, draw=MOSGold, line width=1.5pt, minimum size=6mm, inner sep=0pt},
  dag edge/.style={draw=MOSRed, line width=2pt}
]

% Language World
\node[white, font=\Huge\bfseries] at (0,4.5) {Language};
\node[MOSGold, font=\Large\ttfamily] at (0,3.5) {"Prove..."};

% The Boundary
\fill[MOSRed] (-4, 2.4) rectangle (4, 2.6);

% Morphisms
\draw[-{Latex[length=4mm]}, MOSGold, line width=3pt] (-1, 2.4) -- (-1, 0.5) node[midway, left=2mm] {\Large $\pi$};
\draw[{Latex[length=4mm]}-, MOSGold, line width=3pt] (1, 2.4) -- (1, 0.5) node[midway, right=2mm] {\Large $G$};

% DAG World
\node[white, font=\Huge\bfseries] at (0,-1) {Reasoning DAG};

\node[dag node] (n1) at (-1.5,-2.5) {};
\node[dag node] (n2) at (0,-2.5) {};
\node[dag node] (n3) at (1.5,-2.5) {};
\node[dag node] (n4) at (-1.5,-4) {};
\node[dag node] (n5) at (0,-4) {};

\draw[dag edge] (n1) -- (n2);
\draw[dag edge] (n2) -- (n3);
\draw[dag edge] (n1) -- (n4);
\draw[dag edge] (n2) -- (n5);
\draw[dag edge] (n4) -- (n5);

% Kernel
\node[MOSGold, font=\Large] at (0,-5.5) {$\ker(\pi)$};
\node[white, font=\large] at (0,-6.2) {Semantically void input};

\end{tikzpicture}

\end{center}

%%%%%%%%%%%%%%%%%%%%%%%%%%%%%%%%%%%%%%%%%%%%%%%%%%%%%%%%%%%%%%%%%%%%%%%%%%%%%%
%% VALIDATION
%%%%%%%%%%%%%%%%%%%%%%%%%%%%%%%%%%%%%%%%%%%%%%%%%%%%%%%%%%%%%%%%%%%%%%%%%%%%%%

\begin{highlightbox}

\SectionTitle{Boundary Validation}

\BodyFont

The lift operator is guarded by a validator before a reasoning program
enters the computational engine.

Malformed reasoning DAGs are rejected if they contain

\begin{itemize}

\item duplicate node identifiers,

\item dangling child references,

\item directed cycles,

\item disconnected components,

\item no terminal output node.

\end{itemize}

Only validated reasoning programs are admitted to downstream execution.

\end{highlightbox}

%%%%%%%%%%%%%%%%%%%%%%%%%%%%%%%%%%%%%%%%%%%%%%%%%%%%%%%%%%%%%%%%%%%%%%%%%%%%%%
%% KERNEL
%%%%%%%%%%%%%%%%%%%%%%%%%%%%%%%%%%%%%%%%%%%%%%%%%%%%%%%%%%%%%%%%%%%%%%%%%%%%%%

\begin{mosbox}{Kernel Behaviour}

\BodyFont

Empirically,

semantically irrelevant content is mapped into

\[
\boxed{
\ker(\pi).
}
\]

For example,

non-reasoning text embedded inside a mathematical request generates no
reasoning operators and therefore contributes nothing to the internal
reasoning DAG.

This observed behaviour motivates the interpretation of semantically
void input as belonging to the kernel of the lift.

\end{mosbox}

%%%%%%%%%%%%%%%%%%%%%%%%%%%%%%%%%%%%%%%%%%%%%%%%%%%%%%%%%%%%%%%%%%%%%%%%%%%%%%
%% HONESTY
%%%%%%%%%%%%%%%%%%%%%%%%%%%%%%%%%%%%%%%%%%%%%%%%%%%%%%%%%%%%%%%%%%%%%%%%%%%%%%

\begin{remarkbox}

\BodyFont

\textbf{Claims and Non-Claims}

The notation

\[
L\dashv M
\]

is presented as the intended categorical organisation of the framework.

The corresponding unit and counit identities have \emph{not} been
formally established in the present work.

Accordingly, the adjunction should be interpreted as a guiding
mathematical principle together with the concrete morphisms

\[
(\pi,G),
\]

rather than as a proved categorical theorem.

\end{remarkbox}